% =====================================================================
% Section 3: the case forall x phi (key: univ). Source: research/tracks/universal/notes-final.md
% (and its referee report), with E1/E3(a) from the experiments track. Proofs in app-universal.tex.
% =====================================================================

\section{The case \texorpdfstring{$\forall x\varphi$}{forall x phi}: what instances support}\label{sec:univ}

H\"anni's question starts from statements $\varphi(t)$ and the law $\forall x\varphi$, which ``implies all the other statements'' (\Cref{sec:intro}). In the model of \Cref{sec:model} the data favour the generator that emits exactly the instances. In every comparison analysed (\Cref{thm:univ:B}) they never favour $\forall x\varphi$ over its instance schema $\sigma_\varphi$, except through non-logical weight effects under a selection-aware likelihood (\Cref{prop:univ:B2}); with $\wedge$-rules on quantified formulas the question is open. ``All future data are instances'' is confirmed; ``the axioms prove $\forall x\varphi$'' is not, a Bayesian form of the $\omega$-gap of AS~\S4.5: all closed instances $\varphi(0),\varphi(S0),\dots$ together still do not entail $\forall x\varphi$.

So the intuition is right against hypotheses that fit worse (over-general templates, memorisers), for prediction, and for truth under a Gaifman-type credence, safely in $\N$. It needs three refinements: $\sigma_\varphi$ is a second simple law that implies exactly the data, brought into the class by templates; ``the axioms prove $\forall x\varphi$'' gains from closed instances only through non-logical differences between the laws, as under $\Lsel$ with a background (\Cref{prop:univ:B2}), and with data from a model generator it tends to the prior share of provers in the generator class (\Cref{thm:univ:omega}); and a derivation likelihood rewards the most specific generator that fits, and frequently used lemmas, not the logically strongest law. $\forall x\varphi$ is licensed by an instance at a parameter, a quantified consequence, or a selection-aware likelihood without a guard (\Cref{prop:univ:open}(a), \Cref{prop:univ:quant}(a), \Cref{prop:univ:noguard}(c)), soundly in $\N$ and unsoundly in $\R$; data at parameters are common in proof-assistant practice. Proofs and further tables: \Cref{app:univ}; logarithms are natural.

% ---------------------------------------------------------------------
\subsection{Setup}\label{sec:univ:setup}

\begin{definition}[Instances, schema, term laws]\label{def:univ:setup}\src{universal \S1.1--1.3}
$\varphi(x)$ is an $L_A$-formula with exactly one free variable $x$, which occurs in it; the examples ($0+x=x$, $x+0=x$, $x\cdot0=0$, $\neg Sx=0$, $x<Sx$) are quantifier-free unless stated. $\Ic(\varphi):=\{\varphi(t):t\text{ closed}\}$; $\Io(\varphi)$ also allows parameters in $t$. The instance schema is $\sigma_\varphi:=\varphi[z/x]$, $z$ a $0$-ary term metavariable, closed guard unless stated. Term laws: $\Qnum(S^j0)=(1-q)q^j$, $q=\tfrac12$ by default; $\QGW$, each node picking its root $0,S,+,\cdot$ with probabilities $.4,.3,.2,.1$; $\Qopen:=\rho[p]+(1-\rho)\Qnum$ for a parameter $p$. The first two are root-generated: $\Qg(f(t_1,\dots,t_r))=p_f\prod_i\Qg(t_i)$. One law supplies all bodies and $\forall$E terms.
\end{definition}

The hypotheses\src{universal \S1.3} (\Cref{tab:univ:hyp} in \Cref{app:univ}) are $\Hall=\{\forall x\varphi\}$; $\Hsch=\{\sigma_\varphi\}$; $\Hopen$, $\sigma_\varphi$ with an open guard ($z\sim\Qopen$); $\Hboth=\{\forall x\varphi,\sigma_\varphi\}$ with weight $w_\forall$ on $\forall x\varphi$; memorisers, finite sets of closed instances; over-specific templates such as $\sigma_\varphi[z:=Sz]$; over-general ones, with a subterm or subformula of $\sigma_\varphi$ made a fresh metavariable (down to the inconsistent bare $F_0$); and root splits $\{\sigma_\varphi[z:=f(z_1,\dots,z_r)]:f\text{ a root of }\Qg\}$. For $\varphi=0+x=x$, $\Hall$, $\Hopen$ and $\Hboth$ prove $\forall x\varphi$; $\Hsch$, memorisers, over-specific templates and root splits do not, even over $\Q$ (AS Ex~4.9). $B\oplus_wA$ adds $A$ with weight $w$ to $B$, scaling $B$'s weights by $1-w$. Calculi and likelihoods are as in \Cref{def:model:calculi,def:model:variants}. In $\Cmin$ each chain node cites (probability $1-c$) or applies $\forall$E with $t\sim\Qg$ (probability $c$); $\Copen(c,g)$ adds Gen on the parameter ($g$; cite $p_{\mathrm{cite}}=1-c-g$); $\mu_T$ is the unnormalised derivation mass and $\Lone=\mu_T/Z_T$. H\"anni's scores are $\Sprove$, $\Snc$ and $\Sg$ with $g\equiv1$ on provable data (\Cref{def:model:scores}); for $\beta>1$, $S_{\mathrm{nc},\beta}:=\beta^n\Sg$ with $g_\infty=1/\beta$ rewards each proved datum by a factor $\beta$, and $S_{\mathrm{nc},1}:=\Snc$. Theorems about countable classes assume a proper prior (\Cref{def:model:prior}); the computations use finite or summable families with the prior $2^{-(\sum_A\lvert A\rvert+\#\text{axioms})}$. Defaults: $c=0.3$; $(c,g,\rho,q)=(0.3,0.2,0.1,\tfrac12)$ in $\Copen$; tables state $\alpha$.

% ---------------------------------------------------------------------
\subsection{Instance data do not favour \texorpdfstring{$\forall x\varphi$}{forall x phi} over its schema}\label{sec:univ:never}

\begin{proposition}[Factorisation]\label{prop:univ:factor}\status{proved; computed}\src{universal Prop U1}
In $\Cmin$, let $\varphi$ have one free variable and $m$ leading universal quantifiers. Then $\mu_{\Hall}(s)=c\,\mu_{\Hsch}(s)$ for every sentence $s\neq\forall x\varphi$; $\mu_{\Hall}(\forall x\varphi)=1-c$, $\mu_{\Hsch}(\forall x\varphi)=0$; $Z_{\Hsch}=(1-c)\sum_{k\le m}c^k$, $Z_{\Hall}=(1-c)+cZ_{\Hsch}$. With a background, $\mu_{B\oplus_w\forall x\varphi}=(1-w)\mu_B+wc\,\mu_{\Hsch}$ and $\mu_{B\oplus_w\sigma_\varphi}=(1-w)\mu_B+w\,\mu_{\Hsch}$ on $s\neq\forall x\varphi$. For $k$ variables the factor is $c^k$ on each full instance.
\end{proposition}

\begin{theorem}[Odds on instance data]\label{thm:univ:odds}\status{proved; computed}\src{universal Thm U2}
In $\Cmin$, after $n$ closed instances the odds $O_n:=\pi_n(\Hall)/\pi_n(\Hsch)$ satisfy: $O_n/O_0$ is $1$ under $\Lcl$, under $\Lsel(S)$ with $\forall x\varphi\notin S$, and under the scores ($\forall x\varphi$ consistent); $c^n$ under $\mu_T$ and $\Lmax$; $\big(cZ_{\Hsch}/((1-c)+cZ_{\Hsch})\big)^n$ under $\Lone$, which is $(c/(1+c))^n$ for quantifier-free $\varphi$; $0$ for $n\ge1$ under strict $\Lzero$; and, with a background $B$ under $\mu_T$, $\prod_ir(X_i)$ with $r(s)\in[c,1]$.
\end{theorem}

So at $c=0.3$, under $\Lone$ with quantifier-free $\varphi$, a generator whose axiom is $\forall x\varphi$ outputs $\forall x\varphi$ itself with probability $1/(1+c)=0.77$, and each observed instance multiplies its odds against $\sigma_\varphi$ by $c/(1+c)=0.23$ (\Cref{tab:univ:c2}; checks in \Cref{tab:univ:odds}).

\begin{table}[t]
\centering\footnotesize
\setlength{\tabcolsep}{4pt}
\begin{tabular}{@{}llccccc@{}}
\toprule
likelihood & hypothesis & $n=0$ & $1$ & $5$ & $20$ & $200$\\
\midrule
$\Lcl$ ($=\Lsel$ here) & $\Hall$; $\Hsch$ & .021; .042 & .194; .387 & .303; .605 & .332; .665 & .332; .665\\
 & over-general; over-specific & .926; .011 & .307; .110 & .022; .067 & .000; .000 & .000; .000\\
 & root split, weights $p_f$ & .000 & .002 & .002 & .003 & .003\\
$\mu_T$ & $\Hall$; $\Hsch$ & .021; .042 & .067; .448 & .001; .891 & .000; .996 & .000; .996\\
 & root split, weights $p_f$ & .000 & .002 & .003 & .004 & .004\\
$\Lone$ & $\Hall$; $\Hsch$ & .021; .042 & .053; .455 & .000; .892 & .000; .996 & .000; .996\\
\bottomrule
\end{tabular}
\caption{Running example (c2): mean posterior over 50 seeds, $\varphi=0+x=x$, $\Qnum$, $\Cmin$, prior $2^{-(\text{symbols}+1\text{ per axiom})}$. Class: $\Hall$, $\Hsch$, four over-general templates (among them the inconsistent bare $F_0$, $0.674$ at $n=0$), $\varphi(Sz)$, and the root split with fixed and with learned ($\alpha=1$) weights; the fixed-weight split has exactly $\Hsch$'s law (\Cref{prop:ident:splits}). Four other formulas and $\QGW$ behave alike.}
\label{tab:univ:c2}
\end{table}

\begin{theorem}[Instance data do not favour $\forall x\varphi$]\label{thm:univ:B}\status{proved}\src{universal Thm B}
On data not containing $\forall x\varphi$, the posterior odds of the $\forall$-version against the schema version are at most the prior odds: \textup{(B1)} in $\Cmin$, single axioms, closed instances, under $\Lcl$, $\mu_T$, $\Lone$, $\Lsel$, $\Lmax$ and the scores (equality under $\Lcl$, under $\Lsel$ when $\forall x\varphi\notin S$, and under the scores when $\forall x\varphi$ is consistent); \textup{(B2)} in $\Cmin$, $B\oplus_w\forall x\varphi$ against $B\oplus_w\sigma_\varphi$ for any background $B$ and any data, under $\mu_T$ and $\Lone$, with a fixed $w$ or a common weight prior (equality under $\Lcl$); \textup{(B3)} in $\CUiii$ calculi, quantifier-free $\varphi$, any background (\Cref{prop:univ:sim}), under $\mu_T$ and $\Lmax$; and, for single axioms and quantifier-free data, under the scores, with equality (by Herbrand's theorem, $\{\forall x\varphi\}\vdash s$ iff $\Ic(\varphi)\vdash s$ for quantifier-free $s$).
\end{theorem}

\begin{proof}[Proof idea]
The mechanism is the size principle, through the extra $\forall$E step that $\forall x\varphi$ needs to reach an instance (\Cref{lem:model:size}). (B2) holds at every $w$, and $Z_{\Hall}\ge Z_{\Hsch}$; (B3) rests on a simulation inequality (\Cref{prop:univ:sim} in \Cref{app:univ:odds}).
\end{proof}

Not covered: $\Lsel$ with a background. There $B\oplus_w\forall x\varphi$ has the selected law of $B\oplus_{w_e}\sigma_\varphi$ with the effective weight $w_e=wc/(1-w+wc)$ (proof of \Cref{prop:univ:B2}): with a common fixed $w$ the two theories have different laws and whichever matches the data's mixture proportion wins at a linear rate, and with a learned weight the $\forall$-version can gain a bounded factor (\Cref{prop:univ:B2}). Neither effect is logical, and under the per-citation $\Lselc$ both vanish (\Cref{rem:univ:B2cit}). Also not covered: $\wedge$-detours (\Cref{rem:univ:detour}), normalised likelihoods in $\CUiii$, and the scores with a background or with quantified data, where the $\forall$-version can be favoured (the datum $\neg\exists x\neg\varphi$ is proved by $\Hall$ and not by $\Hsch$).

The direction is robust (\Cref{prop:univ:robust}): in $\Cmin$, at every $s\neq\forall x\varphi$, a full-support noise component leaves the ratio of the laws of $\Hall$ and $\Hsch$ in $[c,1]$ under $\mu_T$, $\le1$ under $\Lone$ and $=1$ under $\Lsel$ with noise after selection; chain weights non-increasing in the number of eliminations keep $P^f_{\Hall}(s)\le P^f_{\Hsch}(s)$; and under $\Lone$, quantifier-free $\varphi$ and any prior on $c$, $O_n/O_0\le2^{-n}$.

\begin{remark}[$\wedge$-detours]\label{rem:univ:detour}\status{refuted (conjecture); counterexample: proof sketch; computed}\src{universal \S2.3}
The factor $c$ of \Cref{prop:univ:sim} depends on the calculus (an earlier version conjectured otherwise; \Cref{app:ver:corrections}): in $\Cand$ at $(p_{\mathrm{cite}},c,a,e)=(0.35,0.15,0.25,0.25)$, $\varphi=0+x=x$, the ratio $\mu_{\Hall}(\varphi(t))/\mu_{\Hsch}(\varphi(t))$ is $0.15553>c$ (\Cref{app:univ:detour}); on a $9139$-point grid it stays $\le0.925$ (normalised $\le0.48$). Whether detours can reverse the direction is open.
\end{remark}

\begin{proposition}[Learned weights under $\Lsel$]\label{prop:univ:B2}\status{proved; computed}\src{universal Prop B2}
In $\Cmin$ let $B=\{\sigma_\psi\}$, $\psi$ quantifier-free, $\inst(\sigma_\psi)\cap\inst(\sigma_\varphi)=\emptyset$, and compare $B\oplus_w\forall x\varphi$ with $B\oplus_w\sigma_\varphi$ (uniform prior on $w$, equal prior masses) under $\Lsel$, $S$ the closed quantifier-free sentences. On data i.i.d.\ from $(1-w^*)\Qg_{\sigma_\psi}+w^*\Qg_{\sigma_\varphi}$ the Bayes factor tends a.s.\ to $h(w^*)$, $h(v):=c/(c+v(1-c))^2\in[c,1/c]$, which exceeds $1$ iff $w^*<\sqrt c/(1+\sqrt c)$.
\end{proposition}

\noindent The cause is the prior induced on the effective weight $w_e$ (computed to $4\cdot10^{-4}$ at $n=10^5$). With learned weights the effect is bounded; it is not evidence that the axioms prove $\forall x\varphi$.

\begin{remark}[Per-citation selection]\label{rem:univ:B2cit}\status{proved}\src{editor, from experiments \S1.7, Prop X6(b) and universal Prop B2}
Under the per-citation $\Lselc$ the selected law of the component $\forall x\varphi$ is $\Qg_{\sigma_\varphi}$, so both theories of \Cref{prop:univ:B2} have law $(1-w)\Qg_{\sigma_\psi}+w\Qg_{\sigma_\varphi}$ at every $w$: the Bayes factor is $1$ (\Cref{rem:model:sel}).
\end{remark}

\noindent A spare-slot prover dies only polynomially: with a uniform prior on $w_\forall$, under $\Lone$ and quantifier-free $\varphi$, $\Hboth$'s mass falls like $\Theta(1/n)$, as $\forall x\varphi$ is a spare slot, a template beside the generator's own that the data never or rarely use (\Cref{prop:univ:both}).

The experiments reproduce \Cref{thm:univ:odds} in their chain $\Cch{1}$ (factor $\tfrac12$: one bit per datum, to $1.1\cdot10^{-11}$) and, under $\Lselc$, an exact tie with $P(T\vdash\forall x\varphi\mid D_n)\to0.484$ ($0.269$ for $\neg Sx=0$), the prior share in their code (E1, \Cref{tab:exp:e1}).

% ---------------------------------------------------------------------
\subsection{Memorisers, over-general and over-specific templates}\label{sec:univ:memo}

Over-general templates lose exponentially (\Cref{thm:univ:size} in \Cref{app:univ:memo}): a law with mass $p<1$ on $\Ic(\varphi)$ loses at least $\ln\frac1p$ nats per datum, a.s.\ and in expectation ($\Ex\,P'(D_n)/P^*(D_n)\le p^n$), and under the citation law ($\Lzero$, $\Lcl$) a first-order template strictly more general than $\sigma_\varphi$, with metavariables i.i.d.\ from $\Qg$ and a formula law $F$, has $p\le\max(\max_fp_f,\max_gF(\mathrm{root}=g))$. Over-specific templates $\sigma_\varphi[z:=f(z_1,\dots,z_r)]$ fail too: under a root-generated $\Qg$ their odds against $\sigma_\varphi$ are a martingale that grows like $p_f^{-n}$ while all data have root $f$ (probability $p_f^n$) and then vanishes (\Cref{prop:univ:overspec}). The memoriser class, however, is not exponentially small in general (\Cref{prop:univ:memo}): with uniform weights and $\Qg$ of infinite support and finite entropy its log-odds against $\Hsch$ fall faster than linearly, but with Laplace weights under $\Qnum$ they stay $\ge-C\ln^2n$ a.s.\ for all large $n$ (conjecturally $-\Theta(\ln^2n)$); under $\QGW$ the computed rate is exponential, as distinct data grow linearly. E1 tests only the single memoriser $\Mem(D_n)$, whose axioms are the distinct data seen so far; every fixed memoriser dies under well-specification (\Cref{rem:ident:memo}).

% ---------------------------------------------------------------------
\subsection{Predictive confirmation}\label{sec:univ:confirm}

\begin{theorem}[Predictive confirmation]\label{thm:univ:confirm}\status{proved; computed}\src{universal Thm U8}
Assume setting W (\Cref{def:ident:W}). Let $M(\cdot\mid D_n)=\sum_T\pi_n(T)P_T$, $G:=\{T:P_T(I)=1\}$, and $T^*\in G$ the generator. \textup{(1)} $\sum_{n\ge0}\Ex[1-M(I\mid D_n)]\le\ln\frac1{\pi(T^*)}$. \textup{(2)} The probability that all future data lie in $I$ is $\pi_n(G)$, and $\pi_n(G)\to1$ a.s. \textup{(3)} $\Ex[1-\pi_n(G)]\le\min\big(1,\sum_{T\notin G}\frac{\pi(T)}{\pi(T^*)}P_T(\operatorname{supp}P_{T^*})^n\big)$. \textup{(4)} If $T\in G$ have law $\hat P$ and $T\notin G$ have law $(1-\varepsilon_T)\hat P+\varepsilon_TN$, $N(I)=0$, then on instance sequences $1-\pi_n(G)=\sum_{T\notin G}\pi(T)(1-\varepsilon_T)^n\big/\big(\pi(G)+\sum_{T\notin G}\pi(T)(1-\varepsilon_T)^n\big)$.
\end{theorem}

\noindent The rate depends on near-universal competitors: $\approx(1-\pi_0)/(\pi_0n)$ for Laplace's rule with a point mass $\pi_0$, but $1-\pi_n(G)=0.37$ at $n=10^{50}$ for a dyadic family ($\varepsilon_j=2^{-j}$ with prior $(1-\pi_0)/(j(j+1))$ on $j\ge1$, $\pi_0=0.01$). Part (2) is the countable-class analogue of Hutter's result for the universal prior, but $G$ contains $\Hsch$, which does not prove $\forall x\varphi$; and confirmation is not stepwise monotone (\Cref{rem:univ:hutter}).

% ---------------------------------------------------------------------
\subsection{The Bayesian \texorpdfstring{$\omega$}{omega}-gap}\label{sec:univ:omega}

\Cref{thm:univ:omega} answers ``do the axioms prove $\forall x\varphi$?''. Its headline is (a) with (b): the posterior probability of a prover tends to the prior share of provers in the generator class $\Cstar:=\{T:P_T=P_{T^*}\}$, which no data move. A likelihood is \emph{faithful} for a calculus $\mathsf C$ if $\operatorname{supp}P_T=\Th_{\mathsf C}(T)\cap X$ for all $T$ ($X$ the data space).

\begin{theorem}[Bayesian $\omega$-gap]\label{thm:univ:omega}\status{(a)--(c), (d1), (e) proved; computed}\src{universal Thm U10}
Let $\Hclass$ be countable with a proper prior $\pi>0$, $\ProvAll:=\{T\in\Hclass:T\vdash\forall x\varphi\}$ (first-order; inconsistent theories included), and the data i.i.d.\ from $P_{T^*}$, $T^*\in\Hclass$, unless stated.
\begin{enumerate}[label=(\alph*),leftmargin=*,itemsep=1pt]
\item $P(T\vdash\forall x\varphi\mid D_n)\to\pi(\ProvAll\cap\Cstar)/\pi(\Cstar)$ a.s.; also for $\vdash_{\mathsf C}$, $\Bel,\Dis,\Indep,\Inc$ (\Cref{def:sound:tri}).
\item \emph{Prior share.} If $\Cstar$ has a prover and a non-prover, the limit is in $(0,1)$: $\Hall,\Hsch$ under $\Lcl$ or $\Lsel$; $B\oplus\forall x\varphi$, $B\oplus\sigma_\varphi$ under $\Lcl$ with a common weight or weight prior. For the two theories of \Cref{prop:univ:B2} under $\Lsel$ the laws differ, but that proposition gives such a limit.
\item \emph{Faithful regime.} (c1) If $\forall x\varphi\in X$, $P(T\vdash_{\mathsf C}\forall x\varphi\mid D_n)\to1[T^*\vdash_{\mathsf C}\forall x\varphi]$. (c2) If also $\mathsf C$ is sound and all axiom instances lie in $X$, members of $\Cstar$ have $T^*$'s first-order theorems and $P(T\vdash\forall x\varphi\mid D_n)\to1[T^*\vdash\forall x\varphi]$. (c3) In $\Cmin$, and $\Copen$ with a guard, $\Lone$ with full-support $\Qg$ is faithful and $\operatorname{supp}P_{\Hsch}=\Ic(\varphi)$ for quantifier-free $\varphi$, so for $T^*=\Hsch$ the limit in (c1) is $0$. If $\mathsf C$ has a rule that maps an instance to a non-instance with positive probability ($\wedge$I, $\vee$I, symmetry of $=$, $\to$I, MP with cited logical axioms), every $T$ with $P_T(\Ic)>0$ emits non-instances: instance-only data are then misspecified for every hypothesis, and the limit is set by the best-fitting generators of the class (\Cref{thm:ident:kl}).
\item \emph{Filtered data} from $\Hall$, seen only in $S=\Ic(\varphi)$. (d1) For quantifier-free $\varphi$ in $\Cmin$ the filtered law is $\Hsch$'s: limit $0$ under $\Lone$, the prior share under $\Lsel$. (d2) \status{proved for the pair; proof sketch in general} For quantified $\varphi$, $\Hsch$ is the KL-minimiser in $\{\Hall,\Hsch\}$ and the limit is $0$; for larger classes the minimiser was not determined.
\item \emph{Unknown filter.} With a countable family $\mathcal F$ of filters $S$ with $P_T(S)>0$ for every $T\in\Hclass$, a prior $\nu$ on $\mathcal F$ independent of $\pi$, and pair laws $P_T(s)1[s\in S]/P_T(S)$, (a) applies to pairs. For $\{\Hall,\Hsch\}$, $\Cmin$, quantifier-free $\varphi$ and data law $\Qg_{\sigma_\varphi}$ the limit is $\frac{\pi(\Hall)\nu(\mathcal F_1)}{\pi(\Hall)\nu(\mathcal F_1)+\pi(\Hsch)\nu(\mathcal F_0)}$, $\mathcal F_0=\{S\supseteq\Ic\}$, $\mathcal F_1=\{S\in\mathcal F_0:\forall x\varphi\notin S\}$.
\end{enumerate}
\end{theorem}

In the prior-share regime the limit lies in $(0,1)$, robustly, but its value depends on the code: for $\varphi=0+x=x$ the share of $\Hall$ in $\{\Hall,\Hsch\}$ is $0.333$ in the symbol-count prior above and in $\pi_1$ of \Cref{def:model:prior}, $0.200$ in $\pi_2$, $0.484$ in the experiments' code and $0.030$--$0.042$ in the per-symbol codes of \Cref{sec:pa} (\Cref{tab:univ:share}).

\noindent H\"anni's own scores leave the question open (\Cref{prop:univ:hanni} in \Cref{app:univ:omega}): for $\varphi=0+x=x$ under $\Sprove$ the posterior tends a.s.\ in total variation to the prior on the consistent provers of $\Ic(\varphi)$, which include independent, true and false theories ($\Q$, $\Q+\forall x\varphi$, $\Q+\neg\forall x\varphi$), so all three trichotomy values keep positive mass and over-generalisation is never penalised; under $\Snc$ and $\Sg$ ($g\equiv1$ on provable data) the provers of all the data share one factor, with the same conclusion, and under $\Snc$ the limit is the prior on the larger set $\{T:T\cup\Ic(\varphi)\text{ consistent}\}$, which contains non-provers such as $\emptyset$. Inconsistent theories inflate $\Bel$ and $\Dis$ alike, so the verifier should use $\Belc$ (\Cref{rem:sound:belc}). Closed instances cannot tell Q4 from its closed schema $z+0=z$, as $(\Q-\text{Q4})$ plus all closed instances of $x+0=x$ does not prove $\forall x(x+0=x)$ (\Cref{prop:univ:q4}). For truth, a credence on structures believes $\forall x\varphi$ in the limit of its true closed instances if $P(\forall x\varphi)>0$ and almost every structure is an elementary extension of its closed-term substructure (Gaifman's condition, cited from memory): safe for $\N$, unsafe for credences that weigh $\R$, $V$, models of $\PA+\neg\Con(\PA)$ or some models of $\Q$ (\Cref{prop:univ:gaifman}).

% ---------------------------------------------------------------------
\subsection{Open instances and the closedness guard}\label{sec:univ:open}

\begin{proposition}[Parameter data and the guard]\label{prop:univ:open}\status{proved; computed}\src{universal Prop U11(a)--(d)}
Assume the closure reading, Gen over parameters, theories without parameters, and a likelihood supported on theorems. \textup{(a)} Once $D$ contains some $\varphi(p)$, every $T$ with $\pi(T\mid D)>0$ proves $\forall x\varphi$. \textup{(b)} If $P_{T^*}(\varphi(p))=\beta>0$, then $P(P(T\vdash\forall x\varphi\mid D_n)<1)\le(1-\beta)^n$. \textup{(c)} In $\Copen(c,g)$ with $\Qopen$, per closed datum under $\Lone$: $\Hopen:\Hsch=\frac{1-\rho}{1+g\rho}$, $\Hall:\Hsch=\frac{c(1-\rho)}{1+c}$, $\Hall:\Hopen=\frac{c(1+g\rho)}{1+c}$. \textup{(d)} If the class is finite and contains the guarded $\Hsch$, closed data give $P(T\vdash\forall x\varphi\mid D_n)\to0$: the guard is learned at rate $\Wo:=\ln\frac{1+g\rho}{1-\rho}$ ($0.1252$ nats per datum at the defaults); the spare-slot prover $\Hboth$ decays only polynomially (like $1/n$ in $\Cmin$, \Cref{prop:univ:both}).
\end{proposition}

\noindent For countable classes with learned-weight provers the limit $0$ is not proved. Here $\rho$ is fixed; in the two-part code of \Cref{sec:pa}, where it is learned, a head-symbol split of $x+0=x$ closes in on the schema at only about 5 bits per doubling of $n$ (\Cref{rem:pa:merge}).

Without a guard ($\DT$, $\lambda$-convention, parameters admissible) the outcome depends on which competitors the class contains; with a selection-aware likelihood the $\omega$-step is taken. The competitors are the unguarded splits $R_k:=\{\varphi(0),\dots,\varphi(S^{k-1}0),\varphi(S^kz)\}$ with weights $(v_0,\dots,v_{k-1},u)$, provers of $\forall x\varphi$ over $\Q$ but not in pure logic; on closed numeral data under $\Lone$ they lose only $q^k\Wo$ nats per datum at their best weights, against $\Wo$ ($\Hopen$) and $\ln\frac{1+c}{c(1-\rho)}$ ($\Hall$; \Cref{prop:univ:rkrate}).

\begin{remark}[Without a guard]\label{rem:univ:openrefuted}\status{refuted; computed}\src{universal \S7.2}
Without a guard the posterior need not take the $\omega$-step on closed data (contrary to an earlier version; \Cref{app:ver:corrections}): in $\{\Hall,\Hopen,R_1\}$, where $R_1=\{\varphi(0),\varphi(Sz)\}$ does not prove $\forall x\varphi$ in pure logic, $P(T\vdash\forall x\varphi\mid D_n)$ is $0.697$ at $n=100$ and $0.000$ from $n=1000$. The claim survives for $\{\Hall,\Hopen\}$ and under $\Lsel$ (cf.\ AS Prop~4.5(b), (c)).
\end{remark}

\begin{proposition}[Who wins without a guard]\label{prop:univ:noguard}\status{proved; computed}\src{universal Prop N3}
Closed data i.i.d.\ from $\Qnum$; the class contains $\Hall$, $\Hopen$, $R_k$ for $k$ in a nonempty $\mathcal K$ (learned or optimal weights), possibly memorisers. In pure logic $\Hall$, $\Hopen$ prove $\forall x\varphi$, and $R_k$ and memorisers do not; over $\Q$ all but the memorisers do. \textup{(a)} $\Lone$, pure logic, provers only $\Hall,\Hopen$: $P(T\vdash\forall x\varphi\mid D_n)\to0$ a.s., exponentially, at rate $\ge(1-q^k)\Wo$ for each $k\in\mathcal K$. \textup{(b)} $\Lone$, over $\Q$, $\mathcal K$ finite: $\to0$ with Laplace-weighted memorisers, $=1$ at every $n$ without them. \textup{(c)} \status{proof sketch} $\Lsel$, finite $\mathcal K$: $\to1$ in pure logic. \textup{(d)} \status{computed; limit conjecture} $\Lone$, $\mathcal K=\{1,\dots,40\}$, memorisers: $1.000$ over $\Q$ from $n=10$ to $10^{12}$, with ever deeper splits.
\end{proposition}

\noindent (b) is proved for Laplace-weighted memorisers; with the Dirichlet($\tfrac12$) weights of \Cref{tab:univ:noguard} the over-$\Q$ value also falls to $0$ with $R_{1..4}$. By (d), a likelihood that ignores the selection approximates the missing guard by ever deeper splits, whose waste $q^k\Wo$ tends to $0$; over $\Q$ the answer is a race between the split family (provers) and the memoriser family (non-provers), a property of the class. A misspecified term law causes a related drift (\Cref{rem:ident:e3a}).

% ---------------------------------------------------------------------
\subsection{Quantified data, lemmas and the equational fragment}\label{sec:univ:quant}

Quantified data decide (\Cref{prop:univ:quant} in \Cref{app:univ:quant}). A datum $s$ with $B\cup\Ic(\varphi)\nvdash s$ gives $B\oplus\sigma_\varphi$ posterior $0$ under any likelihood supported on theorems; over $\Q$, one datum $\forall y(0+Sy=Sy)$ leaves only provers of $\forall x(0+x=x)$. In $\Cmin$, between provers that differ in cost ($\forall x\varphi$ behind $L$ vacuous quantifiers against $\forall x\varphi$, both beside $\sigma_\varphi$, uniform weight priors, equal prior masses) the log Bayes factor tends to $\ln(1+c+\dots+c^L)$ on instance data and grows linearly, a little above $f_qL\ln\frac1c$, when a fraction $f_q>0$ of the data is $\forall x\varphi$ (proof sketch). In $\Lmax$ form a derived sentence becomes an axiom when the derivation steps it saves, at $\ge\ln\frac1c$ each, exceed its prior cost and an $O(\ln n)$ dilution (proof sketch): frequently used lemmas become axioms. So the posterior identifies the data producer's working set of primitives, not a canonical axiomatisation (compare \Cref{prop:pa:usage}). On the equational fragment $E=\{\text{Q4},\text{Q5}\}$ every derivation of $0+S^k0=S^k0$ has at least $k+1$ axiom leaves, so the two-part ($\Lmax$) comparison favours the schema $0+z=z$ linearly in $k$ at $q=\tfrac12$, and the full sum does so at least when $q>0.61$ (\Cref{prop:univ:eqfrag}); full $\Q$, and the full sum at $q=\tfrac12$, are open.

% ---------------------------------------------------------------------
\subsection{Without templates}\label{sec:univ:sentences}

Without templates $\sigma_\varphi$ is unavailable and hypotheses are finite sets of sentences. Setting: $\Cmin$, $c=0.3$, numeral data, $q=\tfrac12$, $\varphi=0+x=x$, learned weights; besides $\Hall$ and memorisers the class has $\Split_k:=\{\varphi(0),\dots,\varphi(S^{k-1}0),\forall y\varphi(S^ky)\}$ and Both0 $:=\{\varphi(0),\forall x\varphi\}$.

Under $\Lcl$, $\Hall$ has the true law and the probability of a prover tends to $1$ (proof sketch). Under $\Lone$, $\Split_k$ loses only $q^k\ln\frac{1+c}c$ nats per datum at its best weights, against $\ln\frac{1+c}c=1.466$ for $\Hall$; it proves $\forall x\varphi$ over $\Q$ but not in pure logic, and memorisers prove it in neither. So in pure logic $P(T\vdash\forall x\varphi\mid D_n)\to0$ in every class computed, and over $\Q$ the limit is a race between families: with depth $\le4$ memorisers win (limit $0$), with depth $\le40$ splits win to $n=10^{12}$ under four prior variants, and with all depths the limit is conjecturally $1$ (\Cref{prop:univ:sentences} in \Cref{app:univ:sentences}). Every finite sentence set that emits infinitely many numeral instances (numeral $\forall$E terms) proves $\forall x\varphi$ over $\Q$ (\Cref{lem:univ:survivors}). $\Split_1$ shows that ``all instances have short derivations, so $\forall x\varphi$ is derivable'', a naive form of Kreisel's conjecture, fails without axioms such as Q3 (\Cref{rem:univ:kreisel}).

% ---------------------------------------------------------------------
\subsection{A good version for this case, and its soundness}\label{sec:univ:intuition}

For $\forall x\varphi$ the good version of \Cref{sec:disc:good} models the selection of reported data by a known filter ($\Lsel$, or the per-citation $\Lselc$, under which the background effect of \Cref{prop:univ:B2} vanishes, \Cref{rem:univ:B2cit}) or by a prior over filters (\Cref{thm:univ:omega}(e)), and its verifier accepts $s$ iff $\Belc(s)\ge1-\delta$. On closed instances it accepts $\forall x\varphi$ only if, eventually, the prior share of provers over theories and filters is $\ge1-\delta$ (for finite $n$ see \Cref{cor:univ:sound}); one parameter instance, or one quantified datum entailing it over the background, suffices.

\begin{corollary}[Soundness for $\forall x\varphi$]\label{cor:univ:sound}\status{proved}\src{universal \S11}
Let the data be i.i.d.\ from $P_{T^*}$ with \emph{fixed} weights, $T^*\in\Hclass$, $T^*\nvdash\forall x\varphi$, and $C^*_\forall:=\{T\in\Cstar:T\vdash\forall x\varphi\Leftrightarrow T^*\vdash\forall x\varphi\}$. The probability that the verifier (on $\Bel$ or $\Belc$) ever accepts $\forall x\varphi$ is at most $\delta/\pi(C^*_\forall)\le\delta/\pi(T^*)$, uniformly in time: \Cref{thm:sound:fixed} read at the sentence $\forall x\varphi$ (proof in \Cref{app:univ:omega}).
\end{corollary}

Without well-specification there is no guarantee. In the misspecified settings of \Cref{thm:univ:omega}(d) and \Cref{prop:univ:noguard,prop:univ:sentences} the limit theories ($\sigma_\varphi$, $R_k$, $\Split_k$, memorisers) are true in $\N$, so in pure logic the verifier errs towards incompleteness; misspecification can also make it accept a falsehood with probability $1$ (\Cref{ex:ident:escape,prop:sound:misspec}). The learned-weight computations here are outside the corollary (see \Cref{thm:sound:avg,thm:sound:shrink}).

\begin{remark}[Splits of $\forall x\varphi$ and the MDL finding]\label{rem:univ:mdl}\status{proved (from the cited results)}\src{universal \S10}
Well specified, a root split with weights $p_f$ ties with the schema under $\Lzero$ and $\Lone$ (\Cref{prop:ident:splits}) and with learned weights pays a polynomial Occam factor (\Cref{prop:ident:splitlzero}). An unmodelled selection creates linear split wins, at rates $(1-q^k)\Wo$ for $R_k$ and $(1-q^k)\ln\frac{1+c}c$ for $\Split_k$ (AS Prop~E.16, split by numeral depth); modelling the selection removes them (\Cref{prop:univ:noguard}(c)). This is \Cref{rem:ident:mdl} for $\forall x\varphi$.
\end{remark}
